% Source: ME17_v1.html

\documentclass[11pt]{article}
\usepackage[margin=1in]{geometry}
\usepackage{amsmath}
\usepackage{amssymb}
\usepackage{siunitx}
\usepackage{enumitem}

\title{ME17: Temperature and Thermal Expansion}
\author{}
\date{}

\begin{document}
\maketitle

\section*{ME17: Temperature and Thermal Expansion}

\begin{enumerate}[leftmargin=*,label=\textbf{Question \arabic*.}]

\item \hfill \textit{10 pts}

Given a temperature of 201 K, convert to degrees Fahrenheit.

\textit{Correct answer:} -97.87

\item \hfill \textit{10 pts}

A metal has a coefficient of linear expansion of 2.6 × 10$^{-5}$ (C°)$^{-1}$.  Suppose a pipe made of this metal is 1.7 m long. Calculate the temperature increase in C° if its length increases by 7.4 mm.

\textit{Correct answer:} 167.4208

\item \hfill \textit{10 pts}

A metal has a coefficient of linear expansion of 1.1 × 10$^{-5}$ (C°)$^{-1}$. Suppose a pipe made of this metal is 5.2 m long. Calculate its length increase in mm if it is warmed by 46 C°.



Δ\emph{L}= \rule{2.5cm}{0.4pt} mm

\textit{Correct answer:} 2.6312

\item \hfill \textit{10 pts}

Suppose a glass bottle with volume 0.012 m$^{3}$ if filled to the brim with a fluid of unknown coefficient of volume expansion. When warmed by 7.1 C°, 16 mL spill onto the lab table. If the glass has linear coefficient 1.9 × 10$^{-5}$K$^{-1}$, calculate the unknown volume coefficient of the fluid in (kC°)$^{-1}$.

(Be careful with units and note that 1000 C°  = 1 kC°.)



 $\beta=$  \rule{2.5cm}{0.4pt} (kC°)$^{-1}$

\textit{Correct answer:} 0.2448

\item \hfill \textit{10 pts}

A bottle filled with 1.3-L of water falls from a height of 150.2 m and hits the ground ( $g=9.8$  m/s$^{2}$). If all the mechanical energy loss of the water went into heating the water, find its \rule{2.5cm}{0.4pt} in temperature (in C$^{o}$).

\textit{Correct answer:} 0.3513

\item \hfill \textit{10 pts}

Suppose 3,141 J of energy are transferred to a block of mass 4.3 kg and specific heat 645 J/(kg⋅C°). Calculate the temperature change in Celsius degrees.

\textit{Correct answer:} 1.1325

\item \hfill \textit{10 pts}

Suppose transferring 24,512 J to a substance of mass 2.8 kg raises the temperature by 11 C°. Calculate the specific heat of the substance in J/(kg⋅C°).

\textit{Correct answer:} 795.8442

\item \hfill \textit{10 pts}

Suppose 5.6 kg of chloroform  at 13.5°C is poured into 9.5 kg of propylene glycol at 36.8°C. Calculate the final equilibrium temperature, neglecting any energy processes associated with mixing, given \emph{c}$_{prop}$ = 2500 J/(kg⋅K) and \emph{c}$_{chloroform}$ = 1050 J/(kg⋅K).



\emph{T} = \rule{2.5cm}{0.4pt} °C

\textit{Correct answer:} 32.1762

\item \hfill \textit{10 pts}

An ice cube tray holds 0.29 kg of water at 17°C. How much energy in kJ much be removed from the water to freeze it?\emph{}

[Use: \emph{L}$_{f}$\emph{}=\emph{}335,000 J/kg\emph{; L}$_{v}$\emph{}= 2,260,000 J/kg; \emph{c}= 4190 J/(kg⋅K)\emph{.}]



\emph{Q}= \rule{2.5cm}{0.4pt} kJ

\textit{Correct answer:} 117.807

\item \hfill \textit{10 pts}

How much energy (in kJ) is required to convert 2.69 kg of ice at -20°C to water at 16°C?

[Use: \emph{c}$_{ice}$ = 2100 J/(kg⋅K); \emph{c}$_{water}$ = 4190 J/(kg⋅K); \emph{L}$_{f,water}$ = 334x10$^{3}$ J/kg]

\textit{Correct answer:} 1,191.7776

\item \hfill \textit{10 pts}

Suppose 7.2 kg of ice at 0°C is placed in 11.5 kg of water at 28.4°C. Calculate the final equilibrium temperature in Celsius. \emph{L}$_{f}$ = 3.35 x 10$^{5}$ J/kg; \emph{c} = 4190 J/(kg⋅K).

[Hint: There are two possible outcomes - the ice will be either fully or partially melted. Start by figuring out which occurs in this case.]

\textit{Correct answer:} 0

\item \hfill \textit{10 pts}

A window has dimensions 1.6 m × 2 m × 13.9 mm where the latter quantity is the thickness in millimeters. The temperature outside is -17°C, and inside it is 18°C. If the thermal conductivity of the glass is 0.8  W/(m⋅K), calculate the magnitude of the energy transfer rate through the window in kW.



\emph{H}= \rule{2.5cm}{0.4pt} kW

\textit{Correct answer:} 6.446

\item \hfill \textit{10 pts}

A ball hanging by a wire has a radius of 3.1 m and emissivity 0.63. If the surface temperature of the ball is 92°C and the room is at temperature 28°C, calculate the net rate in kW at which energy is radiated by the ball.



\emph{H}= \rule{2.5cm}{0.4pt} kW

\textit{Correct answer:} 41.155

\item \hfill \textit{10 pts}

A sheet of metal has a hole cut in it. As the metal is heated, the hole

    [ Select ]

  ["becomes smaller", "becomes larger", "remains the same size", ""]

 .

\begin{itemize}
  \item becomes smaller
  \item becomes larger
  \item remains the same size
\end{itemize}

\item \hfill \textit{10 pts}

Rod A has a twice the specific heat of rod B and three times the linear coefficient of expansion. Both rods have the same mass, length, and initial temperature. After the same energy \emph{Q} is absorbed by each rod, the temperature of rod A will be

    [ Select ]

  ["greater than", "less than", "equal to"]

 the temperature of rod B and the length of rod A will be

    [ Select ]

  ["greater", "equal to", "less than"]

 the length of rod B.

\begin{itemize}
  \item greater than
  \item less than
  \item equal to
  \item greater
  \item equal to
  \item less than
\end{itemize}

\item \hfill \textit{10 pts}

An all-metal pan with metal handle sits on a hot stove. After a while, the handle will become

    [ Select ]

  ["hot", "invisible", "cold", "plastic", "huge"]

, a process called

    [ Select ]

  ["conduction", "radiation", "convection", "cloaking", "progress", "growth"]

.

\begin{itemize}
  \item hot
  \item invisible
  \item cold
  \item plastic
  \item huge
  \item conduction
  \item radiation
  \item convection
  \item cloaking
  \item progress
  \item growth
\end{itemize}

\end{enumerate}

\end{document}
